% CHAPTER 1 - CONDENSED
\chapter{INTRODUCTION}

\section{Background}

Solar photovoltaic (PV) technology has emerged as one of the most viable renewable energy solutions worldwide. Global installed PV capacity exceeded 1,600~GW by 2024, with residential systems comprising approximately 30\% of total installations \cite{irena2024}. The declining cost of PV modules---which dropped over 90\% in the past decade---combined with advances in battery storage technology, has made residential solar systems increasingly accessible.

The Kurdistan Region of Iraq (KRI), located at approximately 36.19$^{\circ}$N latitude and 44.01$^{\circ}$E longitude, presents a unique case study for residential solar deployment. Erbil receives an annual average of approximately 9~hours of sunshine per day, with annual global horizontal irradiation (GHI) ranging between 1,800 and 2,390~kWh/m$^{2}$/year \cite{globalsolaratlas2024}. The Kurdistan Regional Government launched the ``Runaki'' (Light) project in late 2024, aiming to provide 24-hour grid electricity with a progressive tiered tariff structure \cite{krgmoe2024}.

Despite these favorable solar resources, residential deployment faces challenges including chronic grid instability (6--12~hour daily outages), heavy reliance on expensive private diesel generators (\$0.32--0.45/kWh), extreme summer temperatures exceeding 50$^{\circ}$C, and the absence of formal net metering frameworks.


\section{Research Gap}

While international literature extensively documents PV system performance in hot arid climates, limited studies address the specific context of the Kurdistan Region, where:
\begin{itemize}
    \item Grid unreliability fundamentally alters the on-grid vs.\ hybrid economic comparison
    \item Diesel generator displacement represents the primary economic driver rather than grid electricity savings
    \item The new Runaki progressive tariff structure creates novel economic incentives for solar self-consumption
\end{itemize}


\section{Research Objectives}

\begin{enumerate}
    \item To design, plan, and install both on-grid and hybrid residential PV systems in Erbil
    \item To collect and analyze field performance data under real operating conditions
    \item To perform a techno-economic comparison of on-grid versus hybrid configurations
    \item To assess the environmental benefits of diesel generator displacement
    \item To evaluate the impact of the Runaki tariff structure on solar PV economics
\end{enumerate}


\section{Significance}

This study provides the first localized, field-data-driven comparison of on-grid and hybrid residential PV systems in Erbil, directly addressing the practical question facing Kurdistan Region homeowners: which system configuration delivers the best return on investment given the region's unique grid reliability challenges and evolving tariff landscape.
